\documentclass[12pt]{article}
\usepackage{amsmath}
\usepackage{amssymb}
\usepackage[margin=1in]{geometry}

\begin{document}

\begin{center}
\textbf{\Large Homework 4} \\
Tashfeen Hatim Said Omran \\
UCF ID: 5862894 \\
Class: STA4321 \\
October 4, 2026
\end{center}

\textbf{Problem 3.6} \\
I know there are $\binom{5}{2} = 10$ possible pairs of balls that can be selected, and each pair is equally likely with a probability of $1/10$. The possible pairs are $(1,2), (1,3), (1,4), (1,5), (2,3), (2,4), (2,5), (3,4), (3,5), (4,5)$.

a) I define $Y$ as the largest of the two sampled numbers. The possible values for $Y$ are 2, 3, 4, and 5.
\begin{align*}
P(Y = 2) &= P(\{(1,2)\}) = \frac{1}{10} \\
P(Y = 3) &= P(\{(1,3), (2,3)\}) = \frac{2}{10} = \frac{1}{5} \\
P(Y = 4) &= P(\{(1,4), (2,4), (3,4)\}) = \frac{3}{10} \\
P(Y = 5) &= P(\{(1,5), (2,5), (3,5), (4,5)\}) = \frac{4}{10} = \frac{2}{5}
\end{align*}

b) I define $X$ as the sum of the two sampled numbers. The possible values for $X$ are 3, 4, 5, 6, 7, 8, and 9.
\begin{align*}
P(X = 3) &= P(\{(1,2)\}) = \frac{1}{10} \\
P(X = 4) &= P(\{(1,3)\}) = \frac{1}{10} \\
P(X = 5) &= P(\{(1,4), (2,3)\}) = \frac{2}{10} = \frac{1}{5} \\
P(X = 6) &= P(\{(1,5), (2,4)\}) = \frac{2}{10} = \frac{1}{5} \\
P(X = 7) &= P(\{(2,5), (3,4)\}) = \frac{2}{10} = \frac{1}{5} \\
P(X = 8) &= P(\{(3,5)\}) = \frac{1}{10} \\
P(X = 9) &= P(\{(4,5)\}) = \frac{1}{10}
\end{align*}

\textbf{Problem 3.23} \\
I define $Y$ as the gain from playing the game. There are 52 cards in a deck. The probability of drawing a jack or queen is $8/52 = 2/13$, the probability of drawing a king or ace is $8/52 = 2/13$, and the probability of drawing any other card is $36/52 = 9/13$. The probability distribution for $Y$ is:
\begin{align*}
P(Y = 15) &= \frac{2}{13} \\
P(Y = 5) &= \frac{2}{13} \\
P(Y = -4) &= \frac{9}{13}
\end{align*}
I calculate the expected gain $E(Y)$:
\[ E(Y) = \sum y P(y) = 15\left(\frac{2}{13}\right) + 5\left(\frac{2}{13}\right) - 4\left(\frac{9}{13}\right) = \frac{30 + 10 - 36}{13} = \frac{4}{13} \approx 0.31 \]
The expected gain is approximately $\$0.31$.

\textbf{Problem 3.30} \\
I am given a discrete random variable $Y$ with mean $\mu$ and variance $\sigma^2$, and $X = Y + 1$.

a) I expect the mean of $X$ to be larger than $\mu$. Because $X$ is formed by adding 1 to every value of $Y$, the entire probability distribution shifts to the right by 1, which increases the center of mass (mean) by 1.

b) I use Theorems 3.3 and 3.5 to find $E(X)$:
\[ E(X) = E(Y + 1) = E(Y) + E(1) = \mu + 1 \]
Yes, this result agrees with my answer to part (a) since $\mu + 1 > \mu$.

c) I expect the variance of $X$ to be equal to $\sigma^2$. Variance measures the spread or dispersion of a distribution. Shifting the entire distribution by adding a constant changes its location but does not alter how spread out the values are relative to the mean.

\textbf{Problem 3.33} \\
I am given a discrete random variable $Y$ with mean $\mu$ and variance $\sigma^2$, and constants $a$ and $b$.

a) I use the theorems for the expected value of a sum and a constant multiplier:
\[ E(aY + b) = E(aY) + E(b) \]
Using the theorem that $E(cY) = cE(Y)$ and $E(c) = c$, I get:
\[ E(aY + b) = aE(Y) + b = a\mu + b \]

b) I use the definition of variance, $V(X) = E[(X - E(X))^2]$, and substitute $X = aY + b$:
\[ V(aY + b) = E[((aY + b) - E(aY + b))^2] \]
I substitute the result from part (a), $E(aY + b) = a\mu + b$:
\[ V(aY + b) = E[((aY + b) - (a\mu + b))^2] = E[(aY - a\mu)^2] \]
I factor out $a^2$ and apply the constant multiplier theorem for expected values:
\[ V(aY + b) = E[a^2(Y - \mu)^2] = a^2E[(Y - \mu)^2] \]
Since $E[(Y - \mu)^2]$ is the definition of $V(Y)$, I conclude:
\[ V(aY + b) = a^2V(Y) = a^2\sigma^2 \]

\textbf{Problem 3.44} \\
I use a binomial distribution where $Y \sim \text{Binomial}(n=5, p)$ represents the number of successful operations.

a) For $p = 0.8$, I find the probability that all 5 are successful:
\[ P(Y = 5) = \binom{5}{5}(0.8)^5(0.2)^0 = 1(0.32768)(1) = 0.32768 \]

b) For $p = 0.6$, I find the probability that exactly 4 are successful:
\[ P(Y = 4) = \binom{5}{4}(0.6)^4(0.4)^1 = 5(0.1296)(0.4) = 0.2592 \]

c) For $p = 0.3$, I find the probability that less than 2 are successful ($Y = 0$ or $Y = 1$):
\[ P(Y < 2) = P(Y = 0) + P(Y = 1) \]
\[ P(Y = 0) = \binom{5}{0}(0.3)^0(0.7)^5 = 1(1)(0.16807) = 0.16807 \]
\[ P(Y = 1) = \binom{5}{1}(0.3)^1(0.7)^4 = 5(0.3)(0.2401) = 0.36015 \]
\[ P(Y < 2) = 0.16807 + 0.36015 = 0.52822 \]

\textbf{Problem 3.71} \\
I am given a geometric random variable $Y$ with probability of success $p$ and probability of failure $q = 1 - p$.

a) The event $Y > a$ means that the first success occurs after the $a$-th trial. This is equivalent to saying that the first $a$ trials all resulted in failures. Since the trials are independent, the probability of $a$ consecutive failures is the product of their individual probabilities:
\[ P(Y > a) = q^a \]

b) I use the definition of conditional probability and the result from part (a):
\[ P(Y > a + b \mid Y > a) = \frac{P(Y > a + b \cap Y > a)}{P(Y > a)} \]
Since $Y > a + b$ is a subset of $Y > a$, the intersection is $Y > a + b$:
\[ P(Y > a + b \mid Y > a) = \frac{P(Y > a + b)}{P(Y > a)} = \frac{q^{a+b}}{q^a} = q^b = P(Y > b) \]
I think this property is called the \textit{memoryless} property because the probability of waiting $b$ additional trials for a success does not depend on the fact that $a$ failures have already occurred. The system forgets the past failures.

c) In light of the assumptions of identical and independent trials, the result in part (b) is obvious because the outcome of past trials does not influence future trials. After observing $a$ failures, the experiment essentially starts over. The probability of needing more than $b$ additional trials is exactly the same as needing more than $b$ trials from the very beginning.

\textbf{Problem 3.73} \\
I use a geometric distribution where $Y$ is the number of audited accounts until the first substantial error is found. The probability of finding a substantial error is $p = 0.9$, and the probability of not finding one is $q = 0.1$.

a) I find the probability that the first substantial error is the third one audited ($Y = 3$):
\[ P(Y = 3) = q^{3-1}p = q^2p = (0.1)^2(0.9) = (0.01)(0.9) = 0.009 \]

b) I find the probability that the first substantial error occurs on or after the third audited account ($Y \geq 3$):
\[ P(Y \geq 3) = P(Y > 2) \]
Using the property $P(Y > a) = q^a$ from Problem 3.71:
\[ P(Y > 2) = q^2 = (0.1)^2 = 0.01 \]

\end{document}